\subsection{Purpose}

The Best Bike Paths (BBP) system is intended to assist bicycle riders in determining and using the best available bike path given a city's context including factors such as path safety, user-friendly design, obstacles etc.
Users can provide information about the bike path by logging trips or manually and automatically reporting obstacles faced during cycling.
The objective of this Design Document is to show the structure of BBP system in order meet these goal, report about architectural organization and design decisions which have been made to ensure functionalities described in Requirement Analysis and Specification Document (RASD).
The aim of this paper is to be a detailed specification for developers and testers in following phases as implementation, integration and testing phase.

\subsection{Scope}

The BBP system collects, administers and displays bike path and cycling information. More specifically, the system allows for the capturing of bike trips, its includes and validating path-related information, and generating optimal bike paths between any two points based on an aggregated user's input.
This Design Document proposes how the software architecture should be, and what components are necessary to implement these functionalities, its interactions and their interfaces.
This is intended to provide a clear reference for developers and testers during the subsequent phases of implementation, integration, and testing.

\subsection{Definitions}

This subsection reports the definitions of the main concepts used throughout this Design Document. Additional definitions introduced during the requirements analysis phase are described in the RASD.

\begin{itemize}
    \item \textbf{Bike Path}: A path suitable for cycling, either because it includes a dedicated bike lane or because car traffic is limited and compatible with bicycle usage.
    \item \textbf{Trip}: A cycling activity recorded by a registered user, consisting of spatial and temporal data collected during movement.
    \item \textbf{Publishable Information}: Information about bike paths or obstacles that has been confirmed by the user and made visible to the community.
\end{itemize}

\subsection{Acronyms}

\begin{itemize}
    \item \textbf{BBP}: Best Bike Paths
    \item \textbf{DD}: Design Document
    \item \textbf{RASD}: Requirement Analysis and Specification Document
    \item \textbf{UI}: User Interface
    \item \textbf{API}: Application Programming Interface
\end{itemize}

\subsection{Abbreviations}

\begin{itemize}
    \item \textbf{GPS}: Global Positioning System
    \item \textbf{DB}: Database
\end{itemize}

\subsection{Revision History}

Version 1.0 : 07/01/2026.

\subsection{Reference Documents}

This subsection lists the documents that have been consulted and referenced during the preparation of this Design Document and that provide the basis for the architectural and design choices described in the following sections.

\begin{itemize}
    \item Requirement Analysis and Specification Document (RASD) - Best Bike Paths
    \item Software Engineering 2 - Assignment R\&DD, A.Y. 2025--2026
    \item Software Engineering 2 course material and lecture slides
\end{itemize}

\subsection{Document Structure}

This document is structured as follows:

\begin{enumerate}
    \item \textbf{Introduction}: provides a brief description of the project, focusing on the motivations and goals to be achieved through its development, as defined in the RASD.
    \item \textbf{Architectural Design}: describes the architecture of the system at different levels of granularity, including the main components, their interactions, and the adopted architectural patterns.
    \item \textbf{User Interface Design}: presents the organization of the user interface of the platform and describes how users interact with the system.
    \item \textbf{Requirements Traceability}: explains how the requirements defined in the RASD are satisfied by the design elements introduced in this document.
    \item \textbf{Implementation, Integration and Test Plan}: provides an overview of the implementation strategy, integration activities, and testing approach adopted for the system.
    \item \textbf{Effort Spent}: reports the effort spent by each group member in the preparation of this document.
    \item \textbf{Bibliography}: contains the references to the documents and tools used to produce this Design Document.
\end{enumerate}

